\documentclass[]{article}

\usepackage{amsmath}
\usepackage{multirow}
\usepackage{natbib}
\usepackage{graphicx} 


\begin{document}

% Auto-generated figures file
% \usepackage{graphicx}

\begin{figure}[h!]
    \centering
    \includegraphics[width=0.5\textwidth]{../input_files/plots/step_2_r_squared_plot_1775415610.png}
    \caption{Explained Variance Ratio (R-squared) of the two-factor model over the 1,000-day simulation period. The relative stability of the R-squared indicates that the model span is sufficient to capture systematic returns, suggesting that the underperformance of the factor-based covariance estimator is not driven by a failure to explain systematic risk.}
    \label{fig:image1}
\end{figure}

\begin{figure}[h!]
    \centering
    \includegraphics[width=0.5\textwidth]{../input_files/plots/step_2_summary_plot_1775415610.png}
    \caption{Comparison of the factor-based and shrinkage covariance estimators across realized portfolio variance (top), covariance matrix condition number (middle), and portfolio turnover (bottom). The factor-based approach exhibits significant numerical instability, evidenced by a high and volatile condition number, which translates into higher realized portfolio variance compared to the robust shrinkage estimator. Despite this performance gap, both methods result in comparable levels of portfolio turnover, suggesting the rebalancing in the factor-based strategy is driven by estimation noise rather than signal.}
    \label{fig:image0}
\end{figure}



\end{document}
                